\section{Methods}

\subsection{Cohorts, harmonization and the no-imputation data layer}

Data were the baseline (visit V0) of three FACE cohorts (Fondation FondaMental
expert centres): bipolar disorder, schizophrenia and major depression
($N=9{,}013$; BP $6{,}252$, SZ $2{,}209$, DR $552$; 21 recruitment sites).
A harmonized data dictionary maps each modelled variable to its per-cohort source
column, a harmonization rule, and per-variable plausibility (``sanity'') bounds;
out-of-range values are set to missing (never clipped, never imputed). Deterministic
skip-logic decoding sets structural zeros only where a gate is explicitly negative
(e.g.\ attempt count $=0$ when ever-attempted $=0$), logged with provenance; this is
not imputation. Each variable carries a likelihood family and a missingness type.
Identifiers (patient, cohort, DSM-5 subtype, visit, site) are never modelled as
indicators; cohort and subtype enter only as covariates, invariance-grouping
variables, or validation labels. Of $201$ usable common variables, $143$ entered as
modelled indicators across the candidate factors; the remainder are covariates or
identifiers.

\subsection{Candidate ontology and soft priors}

Ten candidate constructs (impulsivity, cognitive flexibility, negative symptoms,
anhedonia, metabolism/immunometabolism, sleep/circadian, overall severity, sensory
abnormality, neurodevelopment, suicidality)---later joined by mania and substance once
their indicators were ingested---define a \emph{soft-prior} ontology rather than a fixed
output. Mapping the candidates against the FACE common-variable coverage fixes the
eligible set, and is itself a primary result: a construct with no common indicator is
reported \emph{not testable}, never back-filled with an invented proxy. Each
(indicator, factor) cell is assigned a prior \emph{role}---primary, plausible-cross,
unlikely or forbidden---whose formal prior is given in Eq.~\eqref{eq:softprior} below;
the default is shrinkage, not exclusion, so the soft-prior exploratory
structural-equation model (ESEM)~\citep{asparouhov2009esem,reise2012bifactor} can
confirm, split, merge, downgrade (proxy), reject or declare not-testable any candidate.

\subsection{Generative model and local independence}

Index patients $i=1,\dots,N$, indicators $j=1,\dots,J$, specific factors $k=1,\dots,K$
and one general factor \Gfac{}. We posit that a few latent causes generate the many
measured quantities: each patient carries a factor vector
$f_i=(G_i,D_{i1},\dots,D_{iK})^{\top}\in\mathbb{R}^{K+1}$ with
\begin{equation}
  f_i \;\sim\; \mathcal{N}_{K+1}\!\left(\mathbf{0},\,\Phi\right),
  \qquad \operatorname{diag}(\Phi)=\mathbf{1},
  \label{eq:factors}
\end{equation}
where the factor \emph{correlation} matrix $\Phi$ has unit diagonal (fixing the latent
scale). Each indicator is a linear reflection of the factors plus item-specific noise,
through a single linear predictor
\begin{equation}
  \eta_{ij} \;=\; \alpha_j \;+\; \lambda_{jG}\,G_i \;+\; \sum_{k=1}^{K}\lambda_{jk}\,D_{ik}
  \;+\; \beta_j^{\top} c_i ,
  \label{eq:linpred}
\end{equation}
with item intercept $\alpha_j$, general loading $\lambda_{jG}$, specific loadings
$\lambda_{jk}$ (row $j$ of the loading matrix
$\Lambda=[\,\lambda_G\mid\Lambda_D\,]\in\mathbb{R}^{J\times(K+1)}$), and item-level
covariates $c_i$ (age, sex, education, site, cohort, assessment year) that adjust an
indicator's \emph{mean} but are never factor indicators. Stacking one patient's items,
with the general column separated from the specific block,
\begin{equation}
  x_i=\mu_i+\underbrace{\lambda_G G_i}_{\text{general}}+\underbrace{\Lambda_D D_i}_{\text{specifics}}+\varepsilon_i,
  \qquad \mu_i=\alpha+B c_i,\qquad \varepsilon_i\sim\mathcal{N}(\mathbf 0,\Psi),
  \label{eq:xstack}
\end{equation}
with diagonal residual covariance $\Psi=\operatorname{diag}(\sigma_1^2,\dots,\sigma_J^2)$.

\paragraph{Local independence (the defining assumption).} Given a patient's true latent
scores, the items are conditionally independent,
\begin{equation}
  p(x_i\mid f_i,\Theta)=\prod_{j\in O_i} p(x_{ij}\mid f_i,\Theta),
  \label{eq:localindep}
\end{equation}
i.e.\ the factors are the only reason items covary. In the Gaussian linear model this is
\emph{identical} to ``$\Psi$ is diagonal'': local independence and diagonal residual
covariance are two faces of one assumption---that the factors explain away all shared
covariance. (It is also why filling missing cells is forbidden: an invented value injects
covariance the model would mistake for a factor.) Because $f_i$ is unobserved and the
likelihood must depend on $\Theta$ alone, we integrate it out by the sum rule and multiply
across independent patients,
\begin{equation}
  L(\Theta)=\prod_{i=1}^{N} p(x_i\mid\Theta)
  =\prod_{i=1}^{N}\int\Bigl[\prod_{j\in O_i} p(x_{ij}\mid f_i,\Theta)\Bigr]\,
  \mathcal{N}(f_i;\mathbf 0,\Phi)\,\mathrm{d}f_i .
  \label{eq:fulllik}
\end{equation}
This is the complete first-principles likelihood---a product over patients (independent
units), a product over items (local independence), and an integral that removes the latent
(the sum rule); everything that follows is the evaluation of \eqref{eq:fulllik}.

\paragraph{Bifactor / ESEM.} Equation~\eqref{eq:linpred} is a bifactor exploratory
structural-equation model~\citep{reise2012bifactor,asparouhov2009esem}: \emph{bifactor}
because every indicator may load on $G$ \emph{and} on its home specific factor;
\emph{ESEM} because an indicator is additionally \emph{allowed} small cross-loadings,
shrunk toward zero by the soft priors \eqref{eq:softprior}---unlike confirmatory factor
analysis, which forbids cross-loadings outright. Two identifications of the factor
covariance are estimated. In the \emph{bifactor} model (primary) $G\perp D_k$ and
$D_k\perp D_l$, so $\Phi$ is block-diagonal, $\Phi=\operatorname{blockdiag}(1,\Phi_D)$;
this is the cleanest first fit. In the \emph{correlated-\Gfac{}} model (sensitivity) the
general--specific block is freed (Eq.~\eqref{eq:corrG}), which is what tests whether
``$\Gfac{}\,\orth\,$biology'' is a finding or an artefact of the bifactor constraint.

\subsection{Mixed-likelihood measurement and the observed-data likelihood}

Rather than coerce variables onto a common scale---which would manufacture covariance and
bias the very factors being estimated---each indicator keeps the observation density
appropriate to its type, all sharing the \emph{same} linear predictor \eqref{eq:linpred}:
\begin{align}
  \text{continuous / $z$-score:}\quad & X_{ij}\sim \mathcal{N}(\eta_{ij},\ \sigma_j^2), \\
  \text{skewed laboratory:}\quad & \log X_{ij}\sim \mathcal{N}(\eta_{ij},\ \sigma_j^2), \\
  \text{ordinal ($C_j$ levels):}\quad & \Pr(X_{ij}\le c)=\operatorname{logit}^{-1}(\tau_{jc}-\eta_{ij}),\ \ \tau_{j1}<\dots<\tau_{j,C_j-1}, \\
  \text{binary:}\quad & X_{ij}\sim \mathrm{Bernoulli}\!\left(\operatorname{logit}^{-1}(\eta_{ij})\right), \\
  \text{count:}\quad & X_{ij}\sim \mathrm{NegBinom}\!\left(\mu_{ij}=e^{\eta_{ij}},\ \phi_j\right).
\end{align}
The continuous block is Gaussian (a Student-$t$ likelihood is available as a robustness
option but is not analytically marginalizable, \S\ref{sub:woodbury}); the non-Gaussian
families fix the factor scale through their cutpoints/offset for identification. Writing
$F_j(\cdot)$ for indicator $j$'s density and $O_i\subseteq\{1,\dots,J\}$ for the
\emph{observed} indicators of patient $i$, the observed-data log-likelihood sums over
observed cells only,
\begin{equation}
  \log p(X_{\mathrm{obs}}\mid\Theta)=\sum_{i=1}^{N}\sum_{j\in O_i}\log F_j\!\left(X_{ij}\mid\eta_{ij},\theta_j\right),
  \label{eq:obslik}
\end{equation}
so a missing cell contributes no term and no completed $N\times J$ matrix is ever formed;
deterministic skip-logic decodes structural zeros as \emph{observed} values (not
imputation).

\subsection{Soft-prior loadings, identification and the rotation problem}

Theory enters through the prior on each loading. With a role
$\rho(j,k)\in\{\textsc{primary},\textsc{cross},\textsc{unlikely},\textsc{forbidden}\}$ per
(indicator, factor) cell,
\begin{equation}
  \lambda_{jk}\mid\rho(j,k)\ \sim\
  \begin{cases}
    \mathcal{N}_{+}(0.70,\,0.25^2) & \textsc{primary (home cell)},\\[2pt]
    \mathcal{N}(0,\,0.25^2)        & \textsc{plausible cross},\\[2pt]
    \mathcal{N}(0,\,0.05^2)        & \textsc{unlikely},\\[2pt]
    \delta_{0}                     & \textsc{forbidden},
  \end{cases}
  \label{eq:softprior}
\end{equation}
where $\mathcal{N}_{+}$ truncates to $(0,\infty)$ and $\delta_0$ fixes the loading at zero.
Remaining priors are weakly informative: $\alpha_j\sim\mathcal{N}(0,1.5^2)$;
$\sigma_j=0.05+\mathrm{HalfNormal}(1)$ (the floor guards against a Heywood case);
ordered-normal cutpoints $\tau_{j\cdot}$; $\phi_j\sim\mathrm{Exponential}(1)$; and
$\Phi\sim\mathrm{LKJ}(\eta{=}2)$, whose density $\propto\det(\Phi)^{\eta-1}$ gently favours
weakly correlated factors~\citep{lewandowski2009lkj}. Latent scores use a non-centred
parameterization ($D_{ik}=z_{ik}$, $z_{ik}\sim\mathcal{N}(0,1)$, scale carried by the
loadings) to defuse the hierarchical funnel.

\paragraph{Why the prior is load-bearing: rotation.} A factor model is not identified from
the likelihood alone. For any invertible $M$,
\begin{equation}
  \Lambda^{\star}=\Lambda M,\quad \Phi^{\star}=M^{-1}\Phi M^{-\top},\quad f^{\star}=M^{-1}f
  \ \Longrightarrow\
  \Lambda^{\star}\Phi^{\star}\Lambda^{\star\top}+\Psi=\Lambda\Phi\Lambda^{\top}+\Psi=\Sigma,
  \label{eq:rotation}
\end{equation}
so rotated parameters imply the identical $\Sigma$, hence the identical Gaussian
likelihood, and are indistinguishable by any amount of data. Confirmatory factor analysis
breaks this by fixing loadings to exactly zero; the soft-prior ESEM breaks it with the
prior. The likelihood is rotation-invariant but the prior is \emph{not}: a rotation moves a
loading the prior wanted near zero away from zero, lowering $p(\Lambda)$, so the posterior
$p(\Lambda\mid X)\propto p(X\mid\Lambda)\,p(\Lambda)$ concentrates at the one orientation
where small loadings align with the soft-zero priors and home loadings are positive and
substantial; the positive truncation removes the residual sign-flip. A corollary gives the
robustness check used below: under a flat loading prior the posterior mode equals the
maximum-likelihood (FIML) solution, so a prior-free refit landing on the same loadings shows
the structure is earned from the data, not manufactured by the priors.

\subsection{The covariance identity and the headline estimand}

For the Gaussian block the inner integral of \eqref{eq:fulllik} is closed-form, giving the
identity at the heart of factor analysis. From \eqref{eq:xstack}, using
$\operatorname{Var}(G_i)=1$, $\operatorname{Cov}(D_i)=\Phi_D$,
$\mathbb{E}[\varepsilon_i\varepsilon_i^{\top}]=\Psi$ diagonal, and the two independence
facts that do the work---$\mathbb{E}[G_i D_i^{\top}]=\mathbf 0$ (bifactor orthogonality)
and $(G_i,D_i)\perp\varepsilon_i$---the cross-terms in
$\Sigma=\operatorname{Cov}(x_i)=\mathbb{E}[(x_i-\mu_i)(x_i-\mu_i)^{\top}]$ vanish and the
survivors sum to
\begin{equation}
  \Sigma=\underbrace{\lambda_G\lambda_G^{\top}}_{\text{general: rank one}}
  +\underbrace{\Lambda_D\Phi_D\Lambda_D^{\top}}_{\text{specific factors}}
  +\underbrace{\Psi}_{\text{private noise}}
  =\Lambda\,\Phi\,\Lambda^{\top}+\Psi,\qquad
  \Phi=\begin{pmatrix}1&\mathbf 0^{\top}\\[1pt] \mathbf 0&\Phi_D\end{pmatrix},
  \label{eq:sigma}
\end{equation}
the zero off-diagonal block being precisely the encoding of $G\perp\text{specifics}$. The
general factor leaves a \emph{rank-one} footprint: every pair of items $j,j'$ covaries
through $G$ by exactly $\lambda_{jG}\lambda_{j'G}$---one shared source touching all items.
Because $x_i$ is a linear combination of jointly Gaussian quantities it is itself Gaussian,
$x_i\sim\mathcal{N}(\mu_i,\Sigma)$.

\paragraph{The headline lives in the off-diagonal block.} The correlated-\Gfac{} model
\emph{unzips} the zero block of \eqref{eq:sigma}, freeing
$\phi_{Gd}=\operatorname{Corr}(G,D)$:
\begin{equation}
  \Phi=\begin{pmatrix}1&\phi_{Gd}^{\top}\\[1pt] \phi_{Gd}&\Phi_D\end{pmatrix},\qquad
  \Sigma=\lambda_G\lambda_G^{\top}
  +\underbrace{\lambda_G\phi_{Gd}^{\top}\Lambda_D^{\top}+\Lambda_D\phi_{Gd}\lambda_G^{\top}}_{\text{the terms orthogonality removed}}
  +\Lambda_D\Phi_D\Lambda_D^{\top}+\Psi .
  \label{eq:corrG}
\end{equation}
The vector $\phi_{Gd}$---the correlation of each specific factor with general burden---is
the headline estimand of the enterprise: ``biology is the least severity-entangled
domain'' is the statement that its metabolic and inflammatory entries are near zero
($0.14$ and $0.06$) while cognition and sleep are far larger ($0.39$ and $0.44$). The
bifactor orthogonality is the \emph{null}; the correlated-\Gfac{} arm measures the
departure from it, and the bifactor direct loading $|\lambda_{jG}|$ and the correlation
$\phi_{Gd}$ agree on the ordering---reporting both is what turns ``biology is separable''
from a constraint artefact into a finding.

\subsection{Missing data, Woodbury computation and factor scores}
\label{sub:woodbury}

\paragraph{Missing data is the marginal property.} For patient $i$ with observed
continuous set $C_i\subseteq O_i$, a foundational property of the multivariate normal is
that the marginal is obtained by \emph{deleting} the unobserved coordinates from
$(\mu,\Sigma)$,
\begin{equation}
  X_{i,C_i}\sim\mathcal{N}\!\left(\mu_{i,C_i},\ \Sigma_{C_i}\right),\qquad
  \Sigma_{C_i}=\Lambda_{C_i}\Phi\Lambda_{C_i}^{\top}+\Psi_{C_i},
  \label{eq:marg}
\end{equation}
with $\Lambda_{C_i}$ the observed-row loadings and
$\Psi_{C_i}=\operatorname{diag}(\sigma_j^2)_{j\in C_i}$. ``Integrating out'' the missing
cells is therefore not an extra step but the marginal property itself; summing
$\log\mathcal{N}(\cdot)$ over each patient's observed sub-vector is \emph{exactly} the
full-information maximum-likelihood (FIML) objective for incomplete multivariate-normal
data. Hence the Bayesian marginalized model and FIML optimize the same observed-data
likelihood and differ only by the priors \eqref{eq:softprior}; a standalone FIML estimator
adds no independent evidence, and confirmation is done in-engine.

\paragraph{Woodbury evaluation at scale.} Evaluating \eqref{eq:marg} naively costs
$\mathcal{O}(|C_i|^3)$ per patient. The Woodbury identity and matrix-determinant lemma
reduce it to the $(K{+}1)$-dimensional factor space,
\begin{align}
  \Sigma_{C_i}^{-1} &= \Psi_{C_i}^{-1}-\Psi_{C_i}^{-1}\Lambda_{C_i}
      \bigl(\Phi^{-1}+\Lambda_{C_i}^{\top}\Psi_{C_i}^{-1}\Lambda_{C_i}\bigr)^{-1}
      \Lambda_{C_i}^{\top}\Psi_{C_i}^{-1}, \label{eq:woodbury}\\[2pt]
  \log\!\left|\Sigma_{C_i}\right| &= \log\!\left|\Phi^{-1}+\Lambda_{C_i}^{\top}\Psi_{C_i}^{-1}\Lambda_{C_i}\right|
      +\log|\Phi|+\!\!\sum_{j\in C_i}\!\log\sigma_j^2 ,
\end{align}
an $\mathcal{O}\!\left((K{+}1)^3\right)$ cost; patients sharing a missingness pattern
share the inner $(K{+}1)\times(K{+}1)$ factorization, computed once per pattern. This
removes the per-patient latent funnel for the high-dimensional continuous block and is what
makes full-sample ($N=9{,}013$) estimation tractable on a workstation (Apple M4 Pro,
24~GB, no GPU); the non-Gaussian block retains explicit per-patient latents.

\paragraph{Factor scores.} Each patient's coordinates are the expected a posteriori (EAP)
factor scores from their observed cells; for an all-continuous pattern,
\begin{equation}
  \hat f_i=\mathbb{E}[f_i\mid X_{i,C_i}]=\Phi\Lambda_{C_i}^{\top}\Sigma_{C_i}^{-1}\!\left(X_{i,C_i}-\mu_{i,C_i}\right),
  \quad
  S_i=\Phi-\Phi\Lambda_{C_i}^{\top}\Sigma_{C_i}^{-1}\Lambda_{C_i}\Phi ,
  \label{eq:scores}
\end{equation}
obtained by MCMC when non-Gaussian cells are present. The per-patient posterior covariance
$S_i$ is exported and propagated downstream (the noise term in the stratification mixture,
\eqref{eq:xd}): a patient with one observed indicator for a factor is \emph{not} treated as
equally characterized as one with six.

\subsection{Staged estimation, priors and convergence}

The global posterior was reached by continuation (homotopy): a well-conditioned
sub-model was fit and then deformed toward the full target one source of difficulty
at a time (S1 continuous core at full $N$; S2 ESEM cross-loadings; S3 mixed-likelihood
suicidality and developmental-risk; S4 thin cohort-specific factors; S5 all factors
jointly plus the correlated-\Gfac{} variant), each converged posterior warm-starting
the next. \emph{Only the global S5 fit is interpreted}; the earlier stages are
convergence/identification checkpoints (publishing an intermediate sub-model would be
exactly the error to avoid). Sampling used the No-U-Turn
sampler~\citep{hoffman2014nuts} via PyMC~\citep{abrilpla2023pymc} with the
NumPyro/JAX backend~\citep{phan2019numpyro} on a workstation (Apple M4 Pro, 24~GB; no
GPU). A stage advanced only on $\hat{R}\le1.01$, effective sample size $\ge400$, zero
divergences, no Heywood cases and acceptable posterior-predictive behaviour; the
reported map was certified at the largest sample that mixed, with a cross-seed
stability guard (Tucker congruence $\varphi=0.993$).

\subsection{In-engine confirmation}

Three checks established that the structure is earned, not manufactured. (i) A
prior-free refit with flat loading priors (identification constraints only) reproduced
the loadings and $\Phi$ essentially exactly (Tucker $\varphi=1.00$); a flat-prior MAP
equals the MLE/FIML. (ii) Posterior-predictive checks compared model-implied and
observed pairwise correlations, yielding a Bayesian SRMR $\approx0.07$ for the
continuous block and reproduction of $21/22$ non-Gaussian indicators' endorsement
rates/means within the $90\%$ predictive interval (lone exception: a $90.8\%$-zero
suicide-attempt count item, an item-level not factor-level misfit). (iii) Model
comparison by WAIC/PSIS-LOO~\citep{watanabe2010waic,vehtari2017loo} decisively
preferred the bifactor over unidimensional and correlated-factor alternatives.
Classical fit indices (CFI/TLI, RMSEA) are not reported---they are weak under heavy
missingness and non-normal indicators and add no independent evidence here.

\subsection{Invariance, robustness and patient scoring}

Measurement invariance across BP/SZ/DR was tested by multi-group loadings and
Bayesian differential item functioning; partials are documented, not hidden.
Robustness to cohort imbalance was assessed by leave-one-cohort-out congruence,
diagnosis-balanced subsampling, site cluster-bootstrap and $1/n_{\text{cohort}}$
weighting (minimum Tucker $\varphi=0.958$); imbalance was handled by weighting and
invariance testing rather than by discarding data. Each patient received, per axis, a
posterior mean, SD and highest-density interval, the number of observed indicators,
and a reliability tier (well-characterized / partial / prior-dominated), so that
downstream analyses propagate uncertainty rather than treating all coordinates as
equally measured.

\subsection{Stratification}

The stratification operates on the certified coordinates $\{x_i,S_i\}_{i=1}^{N}$ with
$x_i=\hat f_i\in\mathbb{R}^{D}$ ($D=9$) and known posterior covariance $S_i$ from
\eqref{eq:scores}; uncertainty is carried throughout rather than discarded.

\paragraph{Structure-discovery gate (is there cluster structure at all?).} Before any
clustering we test the cloud's shape with five complementary statistics, each computed
over posterior draws so that measurement noise cannot manufacture structure.
\emph{Cluster tendency} uses the Hopkins statistic
\begin{equation}
  H=\frac{\sum_{i=1}^{m} u_i^{D}}{\sum_{i=1}^{m} u_i^{D}+\sum_{i=1}^{m} w_i^{D}} ,
\end{equation}
where $w_i$ is the nearest-neighbour distance of a sampled real point and $u_i$ that of
a uniform point over the data support; $H\approx\tfrac12$ indicates no clustering.
\emph{Modality} uses Hartigan's dip statistic~\citep{hartigan1985dip}
$\mathrm{dip}(F_n)=\min_{U\in\mathcal{U}}\sup_x|F_n(x)-U(x)|$ ($\mathcal{U}$ the
unimodal CDFs) on PC1 and each axis. \emph{Model selection} uses the gap
statistic~\citep{tibshirani2001gap},
$\mathrm{Gap}(K)=\mathbb{E}^{*}[\log W_K]-\log W_K$ with within-cluster dispersion
$W_K=\sum_{k}\tfrac{1}{2|\mathcal C_k|}\sum_{i,i'\in\mathcal C_k}\lVert x_i-x_{i'}\rVert^2$,
alongside the Gaussian-mixture BIC profile. \emph{Density} uses
HDBSCAN~\citep{campello2013hdbscan} (fraction of points assigned to noise), and
\emph{topology} a Mapper graph~\citep{singh2007mapper}. A continuum verdict
($H\!\to\!\tfrac12$, unimodal dip, $\mathrm{Gap}$ maximized at $K=1$, flat BIC basin,
HDBSCAN all-noise, a single connected Mapper component) makes archetypes the lead view;
a clustered verdict would promote the mixture.

\paragraph{Model A --- measurement-error Gaussian mixture.} Treating $x_i$ as a noisy
reading of a latent position, $x_i\mid\theta_i\sim\mathcal{N}(\theta_i,S_i)$, with
$\theta_i$ drawn from a $K$-component mixture yields the extreme-deconvolution
likelihood~\citep{bovy2011xd}
\begin{equation}
  x_i\ \sim\ \sum_{k=1}^{K}\pi_k\,\mathcal{N}\!\left(m_k,\ \Sigma_k+S_i\right),
  \qquad \pi\sim\mathrm{Dirichlet}\!\left(\tfrac{\alpha}{K}\mathbf 1\right),
  \label{eq:xd}
\end{equation}
in which \emph{every} observation carries its own additive, known covariance $S_i$, so
imprecise coordinates spread their mass over components (diffuse responsibility) and
prior-dominated or cohort-absent axes self-down-weight. The soft decision regions are
the posterior responsibilities
\begin{equation}
  r_{ik}=\frac{\pi_k\,\mathcal{N}\!\left(x_i\mid m_k,\Sigma_k+S_i\right)}
              {\sum_{l=1}^{K}\pi_l\,\mathcal{N}\!\left(x_i\mid m_l,\Sigma_l+S_i\right)} ,
  \label{eq:resp}
\end{equation}
with optional hard label $\arg\max_k r_{ik}$ and assignment-confidence $\max_k r_{ik}$.
A small Dirichlet concentration ($\alpha\!\ll\!1$) makes \eqref{eq:xd} a
sparse/overfitted finite mixture, so the effective number of regions is data-selected;
$K$ is triangulated by BIC, PSIS-LOO~\citep{vehtari2017loo} and cross-seed stability,
and label switching is resolved by a relabeling algorithm.

\paragraph{Model B --- archetypal analysis (lead view under a continuum).} Each patient
is represented as a convex combination of $A$ extreme phenotypes, each itself a convex
combination of patients~\citep{cutler1994archetypal}:
\begin{equation}
  \min_{W,B}\ \sum_{i=1}^{N}\Bigl\lVert x_i-\sum_{a=1}^{A} w_{ia}\,z_a\Bigr\rVert^2
  \quad\text{s.t.}\quad z_a=\sum_{j=1}^{N} b_{ja}\,x_j,\ \
  w_{ia},b_{ja}\ge 0,\ \ \textstyle\sum_a w_{ia}=1,\ \sum_j b_{ja}=1 ,
  \label{eq:aa}
\end{equation}
i.e.\ $X\approx WZ$ with $Z=BX$ and both $W$ and $B$ row/column-stochastic, which forces
the archetypes $z_a$ onto the convex hull of the cloud. The simplex weights $w_i\in\Delta^{A-1}$
are the soft archetype membership (a patient as ``$0.7$ high-inflammatory $+\ 0.3$
cognitive''). Uncertainty is propagated by refitting \eqref{eq:aa} across posterior
draws and by projecting each draw onto the fixed archetypes (simplex-constrained least
squares), giving a posterior over $w_i$; the number of archetypes $A$ is set by the
explained-variance scree, cross-resample stability (Tucker congruence), and
interpretability, and---under a continuum---is a parsimony choice rather than a natural
number.

\paragraph{Validation.} The partition/archetypes are tested on four gates. Variance
explained per axis,
$\eta^2_d=1-\sum_k\sum_{i\in\mathcal C_k}(x_{id}-\bar x_{kd})^2/\sum_i(x_{id}-\bar x_d)^2$,
quantifies whether structure is driven by the specific/biological axes rather than
\Gfac{} (not-just-severity). Concordance with diagnosis uses the adjusted Rand index
$\mathrm{ARI}=(\mathrm{RI}-\mathbb{E}[\mathrm{RI}])/(\max\mathrm{RI}-\mathbb{E}[\mathrm{RI}])$
and Cram\'er's $V=\sqrt{\chi^2/\!\left(N\min(r-1,c-1)\right)}$ against cohort and the
seven DSM-5 subtypes (transdiagnosticity). Stability is the cross-seed
ARI/Tucker congruence, and a missingness-artefact check compares the accuracy of a
classifier predicting membership from the per-axis coverage pattern against the majority
baseline. A free mixture is compared with one constrained to the DSM-5 partition by
information criteria (the ``tighter description than diagnosis'' test).
UMAP~\citep{mcinnes2018umap} and the principal-component embedding (Results
Fig.~\ref{fig:continuum}) are used for visualization only, never as clustering inputs.

\subsection{Temporal coherence}

Follow-up visits were scored on the fixed measurement model (observed cells,
uncertainty propagated). Longitudinal metric invariance was tested by per-visit
loading congruence vs baseline. Trait/state was decomposed with a measurement-error
random-intercept model per axis, plugging the known baseline measurement variance (so
a low-reliability axis cannot masquerade as state) and including visit fixed effects
(so a shared population trajectory is not miscounted as individual instability),
yielding a population ``slide'' and an individual rank-stability intraclass
correlation. A complementary per-patient geometric analysis counted a coordinate move
only when it cleared measurement error, and quantified archetype-identity persistence
($\kappa$ vs the chance rate). Attrition was characterized by logistic
retention$\sim$baseline coordinates; inverse-probability-of-retention weights were
available for sensitivity analysis.

\subsection{Prognosis}

For each two-year outcome we compared nested models---nuisance + baseline outcome +
severity (the reference), then adding DSM-5 subtype, the map, or both---using an
errors-in-variables Bayesian generalized linear model that propagated each
coordinate's posterior SD. Incremental value was the held-out expected
log-predictive density and, for binary endpoints, the change in area under the ROC
curve with decision-curve analysis. Robustness used inverse-probability weighting for
attrition, restriction to reliably measured axes, and a permutation null. ``Predicts''
denotes incremental association under internal validity, not causation; endpoints are
state transitions from repeated functioning/severity scales, not registered events.

\subsection{Treatment}

Treatment exposures (recovered from per-cohort source dictionaries and harmonized to
common drug classes) were analysed by a causal pipeline: a common-support overlap
gate on the propensity $P(\text{treat}\mid \text{severity}+\text{diagnosis}+
\text{demographics}+\text{the nine coordinates})$~\citep{rosenbaum1983propensity};
covariate balance before/after stabilized inverse-probability-of-treatment
weighting~\citep{austin2011iptw}; and a doubly-robust errors-in-variables moderation
model with the durable-axis $\times$ treatment interaction as the estimand. Average
treatment effects carried an E-value sensitivity analysis for unmeasured
confounding~\citep{vanderweele2017evalue}. A treatment-as-confounder analysis re-fit
the functioning prognosis with and without the drug-class exposures.

\subsection{Reproducibility}

The pipeline is configuration-first with fixed seeds; raw per-cohort data are
confidential and never leave their governed environment, while code, the harmonized
dictionary and regenerated aggregate results are shareable. Every reported number is
reproducible from the numbered pipeline (build $\rightarrow$ fit $\rightarrow$
confirm $\rightarrow$ invariance $\rightarrow$ score $\rightarrow$ robustness
$\rightarrow$ stratify $\rightarrow$ temporal $\rightarrow$ prognosis $\rightarrow$
treatment); Table~\ref{tab:characteristics} is regenerated by an aggregate-only
script.
